

\section{Experiments\label{sec:experiments}}
%We discuss experiments in order to validate our method. 
% \harada{
We test the effectiveness of the proposed semi-supervised ITE estimation method in comparison with various supervised baseline methods, especially when the available labeled data are strictly limited. 
We first conduct experiments using two semi-synthetic datasets based on public real datasets. 
We also design some experiments varying the magnitude of noise on outcomes to explore how the noisy outcomes affect the proposed method.
% Our implementation is available on Github\footnote{\url{https://anonymous.4open.science/r/6ad51d8e-c33c-442f-9a31-f603612bde44/} (Note that this is an anonymous account.)}
%\harada{Our implementation will be available if the paper is accepted.}

% We also conduct experiments using two synthetic datasets. 
% We also conduct experiments using two synthetic datasets. 
%main interests lies in the following three research questions: 

% \vspace{2mm}

% {\bf RQ1:} Does the semi-supervised learning improve the ITE estimation accuracy when the number of labeled instance is small? 

% \vspace{2mm}
% {\bf RQ2:} Is the proposed model robust to the choice of the regularization hyperparameters? 

% \noindent
% {\bf RQ2:} Is the proposed model robust to noisy outcomes?
%\footnote{\harada{Our Implementation of this study is available ..}}
% In both setting, our goal is to explore how the performances change in response to the size of labeled data and validate the advantage of the proposed method.}

% \kashima{Sometimes it is nice to first give research questions we want to answer through experiments. For example, we can say like ``We conducted experiments using .... to investigate the following questions: (1) does the semi-supervised learning improves the estimation accuracy when the number of labeled instance is small, (2) .... "}

\subsection{Datasets}

Owing to the counterfactual nature of ITE estimation, we rarely access real-world datasets including ground truth ITEs, and therefore cannot directly evaluate ITE estimation methods like the standard supervised learning methods using cross-validation.
Therefore, following the existing work~\cite{johansson2016learning}, we employ two semi-synthetics datasets whose counterfactual outcomes are generated through simulations. 
Refer to the original papers for the details on outcome generations~\cite{hill2011bayesian,johansson2016learning}.


\subsubsection{News dataset} is a dataset including opinions of media consumers for news articles~\cite{johansson2016learning}. 
It contains $5{,}000$ news articles and outcomes generated from the NY Times corpus\footnote{\url{https://archive.ics.uci.edu/ml/datasets/Bag+of+Words}}.
Each article is consumed on desktop~($t=0$)~or mobile~($t=1$) and it is assumed that media consumers prefer to read some articles on mobile than desktop. 
% Each article is represented using bag of words and the size of the vocabulary is $3{,}477$. 
Each article is generated by a topic model and represented in the bag-of-words representation. The size of the vocabulary is $3{,}477$. 
% \harada{We use the subset of dataset employed in the previous work}~\cite{johansson2016learning}.


\subsubsection{IHDP dataset} is a dataset created by randomized experiments called the Infant Health and Development Program (IHDP)~\cite{hill2011bayesian} to examine the effect of special child care on future test scores. 
It contains the results of $747$ subjects~($139$ treated subjects and $608$ control subjects)~with $25$ covariates related to infants and their mothers.
Following the existing studies~\cite{johansson2016learning,shalit2017estimating}, the ground-truth counterfactual outcomes are simulated using the NPCI package~\cite{dorie2016npci}. 


\subsubsection{Synthetic dataset} is a synthetically generated dataset for this study. 
We generate $1{,}000$ instances that have eight covariates sampled from $\mathcal{N}({\bf 0}^{8\times1},0.5\times(\Sigma+\Sigma^{\top} ))$,
where $\Sigma $ is sampled from $\sim\mathcal{U}((-1,1)^{8\times8})$.
The treatment $t$ is sampled from $\text{Bern}(\sigma({\bf w}_t^{\top}{\bf x}+\epsilon_t))$, where $\sigma(\cdot)=\frac{1}{1+\exp(-\cdot)}, {\bf w}_t\sim\mathcal{U}((-1,1)^{8\times1}), \epsilon_t\sim\mathcal{N}(0, 0.1)$. 
The treatment outcome and control outcome are generated as  $y^{1}|\x\sim \sin({\bf w}^{\top}_y \x) +c\epsilon_{y}$ and $y^{0}|\x\sim \cos({\bf w}^{\top}_y \x) +c\epsilon_{y}$,
respectively, where ${\bf w} \sim\mathcal{U}((-1,1)^{8\times1})$ and $ \epsilon_y\sim\mathcal{N}({\bf 0}, 1)$. 


% \harada{\kashima{Any reason to use a subset? $\rightarrow$}
% We use the subset of dataset employed in the previous work}~\cite{shalit2017estimating}.\haradap{No. I will delete this sentence.}

% \subsubsection{Synthetic dataset}\harada{write systems}. 




% \subsection{RQ1: Does the semi-supervised learning improve the ITE estimation accuracy when the number of labeled instance is small? }

\subsection{Experimental settings}
% \subsubsection{Experimental settings}
Since we are particularly interested in the situation when the available labeled data are strictly limited, we split the data into a training dataset, validation dataset, and a test dataset by limiting the size of the training data. 
We change the ratio of the training to investigate the performance; we use $10\%, 5\%,$ and $1\%$ of the whole data from the News dataset, and use $40\%, 20\%$, and $10\%$ of those from the IHDP dataset for the training datasets. 
% On IHDP dataset, we hold $50\%$ and $10\%$ of the whole data for test and validation, respectively. 
The rest $80\%$ and $10\%$ of the whole News data are used for test and validation, respectively.
Similarly, $50\%$ and $10\%$ of the whole IHDP dataset are used for test and validation, respectively.
% On the News dataset, we hold $80\%$ and $10\%$ for test and validation. 
% We report the average results of $10$ trials on the News dataset and $50$ trials on the IHDP dataset. 
We report the average results of $10$ trials on the News dataset, $50$ trials on the IHDP dataset, and $10$ trials on the Synthetic datasets.


In addition to the evaluation under labeled data scarcity, we also test the robustness against label noises. 
As pointed out in previous studies, noisy labels in training data can severely deteriorate predictive performance, especially in semi-supervised learning.  
Following the previous  work~\cite{hill2011bayesian,johansson2016learning}, we add the noise $\epsilon\sim\mathcal{N}(0, c^2)$ to the observed outcomes in the training data, where $c \in \{1, 3, 5, 7, 9\}$.
%Note that in both datasets, the outcomes employed in the previous studies~\cite{hill2011bayesian,johansson2016learning,shalit2017estimating} and this study are generated as $c=1$ unless mentioned explicitly.  
In this evaluation, we use $1\%$ of the whole data as the training data for the News dataset and $10\%$  for the IHDP dataset, respectively, since we are mainly interested in label-scarce situations.



The hyper-parameters are tuned based on the prediction loss using the  observed outcomes on the validation data. 
We calculate the similarities between the instances by using the Gaussian kernel; 
we select $\sigma^{2}$ from $\{5\times10^{-3}, 1\times10^{-3},\ldots , 1\times10^{2}, 5\times10^{2}\}$, 
and select $\lambda_o$ and $\lambda_e$ from $\{1\times10^{-3}, 1\times10^{-2}, \ldots, 1\times10^{2}\}$. 
Because the scales of treatment outcomes and control outcomes are not always the same, we found scaling the regularization terms according to them is beneficial; 
% incorporating them in setting of the regularization hyperparameters is beneficial according to our observations;
specifically, we scale the regularization terms with respect to the treatment outcomes, the control outcomes, and the treatment effects by $\alpha={1}/{\sigma_{y^1}^{2}}, \beta={1}/{\sigma_{y^0}^{2}}$, and $\gamma={1}/{(\sigma_{y^1}^{2}+\sigma_{y^0}^{2})}$, respectively.
%, and use them for the regularization hyperparameters \harada{as $\alpha\lambda_o$, $\beta\lambda_o$, $\gamma\lambda_e$ for treatment outcome, control outcome, and treatment effect regularization, respectively}.
We apply principal component analysis to reduce the input dimensions before applying the Gaussian kernel; we select the number of dimensions from $\{2, 4, 6, 8, 16, 32, 64\}$. 
The learning rate is set to $1\times10^{-3}$ and the mini-batch sizes $b_1, b_2$ are chosen from $\{4,8,16,32\}$. 

As the evaluation metrics, we report the {\it Precision in Estimation of Heterogeneous Effect~(PEHE)~} used in the previous research~\cite{hill2011bayesian}.
PEHE is the estimation error of individual  treatment effects, and is defined as 
\[
\epsilon_{\rm PEHE}=\frac{1}{N+M}\sum_{i=1}^{N+M}(\tau_i - \hat{\tau}_i)^2.
\]
Following the previous studies~\cite{shalit2017estimating,yao2018representation}, we evaluate the predictive performance for labeled instances and unlabeled instances separately. 
% The first task is to estimate ITE for instances which are included in the train dataset~(which we refer to as `within-sample'). 
% Although we observe factual outcomes of instances in the training data,  this setting still evaluate predictive performance because we cannot access to the counterfactual outcomes. 
% The other task is to estimate ITE for instances which are not included in the training data~(which we call `without-sample'). 
% The first task is conducted using the training and validation dataset in the same way as existing work~\cite{shalit2017estimating}. 
%The first task is to estimate the ITEs for the labeled instances. 
Note that, although we observe the factual outcomes of the labeled data, their true ITEs are still unknown because we cannot observe their counterfactual outcomes.
%Note that we do not know them because only the factual outcomes are given in the labeled data.
%of instances in the training data,  this setting still evaluate predictive performance because we cannot access to the counterfactual outcomes. 
%The other task is to estimate the ITEs for unlabeled instances.
%The first task is conducted using the training and validation datasets following the procedure proposed in the previous work~\cite{shalit2017estimating}. 

\begin{table*}[tbp]
\caption{The performance comparison of different methods on News dataset. The $^\dagger$ indicates that our proposed method (CP) performs statistically significantly better than the baselines by the paired $t$-test ($p<0.05)$. The bold results indicate the best results in terms of the average.}\label{table:news}
% \scalebox{0.919}[0.92]
{
%   \begin{center}
\centering

%     \begin{tabular}{c}
      % 1
%       \begin{minipage}{\hsize}
\begin{tabular} {ccccccccc }
% \caption{Experimental results}

\toprule[2pt]
          $ \sqrt{\epsilon_{\rm PEHE}} $         & \multicolumn{2}{c}{ \bf{News} 1\%}        &              & \multicolumn{2}{c}{\bf{News} 5\%}& &  \multicolumn{2}{c}{\bf{News} 10\%}      \\[3pt]\cmidrule{2-3}\cmidrule{5-6}\cmidrule{8-9}
Method                      & labeled  & unlabeled & \multicolumn{1}{l}{} & labeled & unlabeled & \multicolumn{1}{l}{} & labeled & unlabeled\\[3pt] \cmidrule{1-9}
Ridge-1   &  $^\dagger4.494_{\pm{1.116}}$  &   $^\dagger4.304_{\pm{0.988}}$         & \multicolumn{1}{l}{} &  $^\dagger4.666_{\pm{1.0578}}$ &  $^\dagger3.951_{\pm{0.954}}$ &   & $^\dagger4.464_{\pm{1.082}}$   & $^\dagger3.607_{\pm{0.943}}$      \\[3pt]
Ridge-2   &  $2.914_{\pm{0.797}}$ &  $^\dagger2.969_{\pm{0.814}}$  & \multicolumn{1}{l}{} & $^\dagger2.519_{\pm{0.586}}$ & $^\dagger2.664_{\pm{0.614}}$ &  & $^\dagger2.560_{\pm{0.558}}$ & $^\dagger2.862_{\pm{0.621}}$      \\[3pt]

Lasso-1   &        $^\dagger4.464_{\pm{1.082}}$   &   $^\dagger3.607_{\pm{0.943}}$   & \multicolumn{1}{l}{} &    $^\dagger4.466_{\pm{1.058}}$    &     $^\dagger3.367_{\pm{0.985}}$    &  &   $^\dagger4.464_{\pm{1.0822}}$& $^\dagger3.330_{\pm{0.984}}$      \\[3pt]
Lasso-2   &        $^\dagger3.344_{\pm{1.022}}$   &   $^\dagger3.476_{\pm{1.038}}$   & \multicolumn{1}{l}{} &    $^\dagger2.568_{\pm{0.714}}$   &  $^\dagger2.848_{\pm{0.751}}$ &  &  $^\dagger2.269_{\pm{0.628}}$  &  $^\dagger2.616_{\pm{0.663}}$      \\[3pt]

\cmidrule{1-9}
\multicolumn{1}{c}{$k$NN} &    $^\dagger3.678_{\pm{1.250}}$   &   $^\dagger3.677_{\pm{1.246}}$  & \multicolumn{1}{l}{} &   $^\dagger3.351_{\pm{1.004}}$&   $^\dagger3.434_{\pm{1.018}}$  &  &  $^\dagger3.130_{\pm{0.752}}$ & $^\dagger3.294_{\pm{0.766}}$      \\[3pt]
\multicolumn{1}{c}{PSM}    &  $^\dagger3.713_{\pm{1.149}}$  &  $^\dagger3.662_{\pm{1.127}}$ & \multicolumn{1}{l}{} &   $^\dagger3.363_{\pm{0.901}}$  &  $^\dagger3.500_{\pm{0.961}}$  &  &   $^\dagger3.260_{\pm{0.734}}$& $^\dagger3.526_{\pm{0.832}}$      \\[3pt]
\cmidrule{1-9}
\multicolumn{1}{c}{RF}    &  $^\dagger4.494_{\pm{1.116}}$  &    $^\dagger3.691_{\pm{0.878}}$   & \multicolumn{1}{l}{} &    $^\dagger4.466_{\pm{1.058}}$   &  $^\dagger2.975_{\pm{0.874}}$    &  &  $^\dagger4.464_{\pm{1.082}}$ & $^\dagger2.657_{\pm{0.682}}$     \\[3pt]
\multicolumn{1}{c}{CF}    &  $^\dagger3.691_{\pm{1.082}}$  &    $^\dagger3.607_{\pm{0.943}}$   & \multicolumn{1}{l}{} &  $^\dagger3.196_{\pm{0.901}}$ & $^\dagger3.215_{\pm{0.910}}$ &  & $^\dagger3.101_{\pm{0.806}}$ & $^\dagger3.129_{\pm{0.818}}$ \\[3pt]
\cmidrule{1-9}
\multicolumn{1}{c}{TARNET}    &        $^\dagger3.166_{\pm{0.742}}$       &       $^\ddagger3.160_{\pm{0.722}}$         & \multicolumn{1}{l}{} &   $^\dagger2.670_{\pm{0.796}}$  &   $^\dagger2.666_{\pm{0.773}}$      &  &   $^\dagger2.589_{\pm{0.894}}$ & $^\dagger2.598_{\pm{0.869}}$      \\[3pt]
\multicolumn{1}{c}{CFR}    &        $2.908_{\pm{0.752}}$       &       $2.925_{\pm{0.746}}$         & \multicolumn{1}{l}{} &   $^\dagger2.590_{\pm{0.772}}$   &  $^\dagger2.546_{\pm{0.796}}$   &  &   $^\dagger2.570_{\pm{0.519}}$& $^\dagger2.451_{\pm{0.547}}$      \\[3pt]
\cmidrule{1-9}
CP (proposed) &        $\bf 2.844_{\pm{0.683}}$       &       $\bf 2.823_{\pm{0.656}}$        & \multicolumn{1}{l}{} &     $\bf 2.310_{\pm{0.430}}$          & $\bf 2.446_{\pm{0.471}}$       &  &   $\bf 2.003_{\pm{0.393}} $& $\bf 2.153_{\pm{0.436}}$      \\[3pt]
% \cmidrule{1-9}
% Improvement\\
\bottomrule[2pt]
    \end{tabular}
    }
\end{table*}
\begin{table*}[tbp]
\caption{The performance comparison of different methods on IHDP dataset. The $^\dagger$ indicates that our proposed method (CP) performs statistically significantly better than the baselines by the paired $t$-test ($p<0.05$). The bold results indicate the best results in terms of average.}\label{table:ihdp}
\scalebox{0.975}[0.975]
{
\centering
\begin{tabular} {ccccccccc }
\toprule[2pt]
       $ \sqrt{\epsilon_{\rm PEHE}} $         & \multicolumn{2}{c}{ \bf{IHDP} 10\%}        &              & \multicolumn{2}{c}{\bf{IHDP} 20\%}& &  \multicolumn{2}{c}{\bf{IHDP} 40\%}      \\[3pt]\cmidrule{2-3}\cmidrule{5-6}\cmidrule{8-9}
Method                      & labeled & unlabeled & \multicolumn{1}{l}{} & labeled & unlabeled & \multicolumn{1}{l}{} & labeled & unlabeled\\[3pt] \cmidrule{1-9}
Ridge-1   &        $^\dagger5.484_{\pm{8.825}}$ &  $^\dagger5.696_{\pm{7.328}}$ & \multicolumn{1}{l}{} &   $^\dagger5.067_{\pm{8.337}}$  &  $^\dagger4.692_{\pm{6.943}}$       &   &$^\dagger4.80_{\pm{8.022}}$   & $^\dagger4.448_{\pm{6.874}}$      \\[3pt]
Ridge-2   &        $^\dagger3.426_{\pm{5.692}}$      &       $^\dagger3.357_{\pm{5.177}}$       & \multicolumn{1}{l}{} &  $^\dagger2.918_{\pm{4.874}}$  &    $^\dagger2.918_{\pm{4.730}}$  &   &  $^\dagger2.605_{\pm{4.314}}$ & $^\dagger2.639_{\pm{4.496}}$      \\[3pt]

Lasso-1   &        $^\dagger6.685_{\pm{10.655}}$       &       $^\dagger6.408_{\pm{9.900}}$         & \multicolumn{1}{l}{} &        $^\dagger6.435_{\pm{10.147}}$       &  $^\dagger6.2446_{\pm{9.639}}$       &     &   $^\dagger6.338_{\pm{9.704}}$ &  $^\dagger6.223_{\pm{9.596}}$  \\[3pt]
Lasso-2   &        $^\dagger3.118_{\pm{5.204}}$       &       $^\ddagger3.292_{\pm{5.725}}$         & \multicolumn{1}{l}{} &         $^\dagger2.684_{\pm{4.428}}$    &  $^\dagger2.789_{\pm{4.731}}$       &  &   $^\dagger2.512_{\pm{4.075}}$& $^\dagger2.571_{\pm{4.379}}$      \\[3pt]
\cmidrule{1-9}
\multicolumn{1}{c}{$k$NN} &        $^\dagger4.457_{\pm{6.957}}$  &  $^\dagger4.603_{\pm{6.629}}$         & \multicolumn{1}{l}{} &      $^\dagger4.023_{\pm{6.193}}$     &    $^\dagger4.370_{\pm{6.244}}$     &  &   $^\dagger3.623_{\pm{5.316}}$ & $^\dagger4.109_{\pm{5.936}}$      \\[3pt]
\multicolumn{1}{c}{PSM}    &        $^\dagger6.506_{\pm{10.077}}$       &       $^\dagger6.982_{\pm{10.672}}$         & \multicolumn{1}{l}{} &   $^\dagger6.277_{\pm{9.708}}$     &    $^\dagger7.209_{\pm{11.077}}$    &  &   $^\dagger6.065_{\pm{9.362}}$ & $^\dagger7.181_{\pm{9.362}}$      \\[3pt]
\cmidrule{1-9}
\multicolumn{1}{c}{RF}    &        $^\dagger6.924_{\pm{10.620}}$       &       $^\dagger5.356_{\pm{8.790}}$         & \multicolumn{1}{l}{} &   $^\dagger6.854_{\pm{10.471}}$      &    $^\dagger4.845_{\pm{8.241}}$     &  &   $ ^\dagger6.928_{\pm{10.396}}$ & $^\dagger4.549_{\pm{7.822}}$\\[3pt]
\multicolumn{1}{c}{CF}    &        $^\dagger5.389_{\pm{8.736}}$       &       $^\dagger5.255_{\pm{8.070}}$         & \multicolumn{1}{l}{} &   $^\dagger4.939_{\pm{7.762}}$      &    $^\dagger4.955_{\pm{7.503}}$   &  &   $^\dagger4.611_{\pm{7.149}}$  &  $^\dagger4.764_{\pm{7.448}}$\\[3pt]
\cmidrule{1-9}
\multicolumn{1}{c}{TARNET}    &        $^\dagger3.827_{\pm{5.315}}$       &       $^\dagger3.664_{\pm{4.888}}$         & \multicolumn{1}{l}{} &   $^\dagger2.770_{\pm{3.617}}$  &    $^\dagger2.770_{\pm{3.542}}$   &     &   $^\dagger2.005_{\pm{2.447}}$& $^\dagger2.267_{\pm{2.825}}$      \\[3pt]
\multicolumn{1}{c}{CFR}    &        $^\dagger3.461_{\pm{5.1444}}$       &       $^\dagger3.292_{\pm{4.619}}$         & \multicolumn{1}{l}{} &  $^\dagger2.381_{\pm{3.126}}$ &    $^\dagger2.403_{\pm{3.080}}$      &  &   $1.572_{\pm{1.937}}$ & $1.815_{\pm{2.204}}	$   	 \\[3pt]
\cmidrule{1-9}
CP~(proposed) &       $\bf 2.427_{\pm{3.189}}$ & $\bf 2.652_{\pm{3.469}}$        & \multicolumn{1}{l}{} &     $\bf 1.686_{\pm{1.838}}$          & $\bf 1.961_{\pm{2.343}}$       &  &   $\bf 1.299_{\pm{1.001}} $& $\bf 1.485_{\pm{1.433}}$      \\[3pt]
\bottomrule[2pt]
    \end{tabular}
    }
\end{table*}

\begin{table*}[tbp]
\caption{The performance comparison of different methods on Synthetic dataset. The $^\dagger$ indicates that our proposed method (CP) performs statistically significantly better than the baselines by the paired $t$-test ($p<0.05$). The bold results indicate the best results in terms of average.}\label{table:synthetic}
\scalebox{0.975}[0.975]
{
\centering
\begin{tabular} {ccccccccc }
\toprule[2pt]
       $ \sqrt{\epsilon_{\rm PEHE}} $         & \multicolumn{2}{c}{ \bf{Synthetic} 10\%}        &              & \multicolumn{2}{c}{\bf{Synthetic} 20\%}& &  \multicolumn{2}{c}{\bf{Synthetic} 40\%}      \\[3pt]\cmidrule{2-3}\cmidrule{5-6}\cmidrule{8-9}
Method                      & labeled & unlabeled & \multicolumn{1}{l}{} & labeled & unlabeled & \multicolumn{1}{l}{} & labeled & unlabeled\\[3pt] \cmidrule{1-9}
Ridge-1   &        $^\dagger0.966_{\pm{0.078}}$ &  $^\dagger0.960_{\pm{0.78}}$ & \multicolumn{1}{l}{} &   $^\dagger0.965_{\pm{0.072}}$  &  $^\dagger0.960_{\pm{0.79}}$       &   & $^\dagger0.959_{\pm{0.062}}$ &  $^\dagger0.949_{\pm{0.77}}$      \\[3pt]
Ridge-2   &        $^\dagger0.890_{\pm{0.180}}$  &  $^\dagger0.950_{\pm{0.193}}$       & \multicolumn{1}{l}{}
&  $^\dagger0.858_{\pm{0.164}}$  &    $^\dagger0.881_{\pm{0.193}}$  &   &  $^\dagger0.836_{\pm{0.137}}$ & $^\dagger0.855_{\pm{0.175}}$      \\[3pt]

Lasso-1   &        $^\dagger0.965_{\pm{0.077}}$       &       $^\dagger0.959_{\pm{0.075}}$         & \multicolumn{1}{l}{} &        $^\dagger0.964_{\pm{0.069}}$       &  $^\dagger0.951_{\pm{0.074}}$       &     &   $^\dagger0.959_{\pm{0.061}}$ &  $^\dagger0.949_{\pm{0.074}}$  \\[3pt]

Lasso-2   &        $^\dagger0.862_{\pm{0.121}}$       &       $^\ddagger0.887_{\pm{0.142}}$         & \multicolumn{1}{l}{} &         $^\dagger0.877_{\pm{0.111}}$    &  $^\dagger0.876_{\pm{0.126}}$       &  &   $^\dagger0.872_{\pm{0.104}}$& $^\dagger0.871_{\pm{0.124}}$     \\[3pt]\cmidrule{1-9}

\multicolumn{1}{c}{$k$NN} &        $^\dagger0.787_{\pm{0.104}}$  &  $^\dagger0.842_{\pm{0.126}}$         & \multicolumn{1}{l}{} &      $^\dagger0.703_{\pm{0.107}}$     &    $^\dagger0.766_{\pm{0.134}}$     &  &   $^\dagger0.643_{\pm{0.096}}$ & $^\dagger0.695_{\pm{5.936}}$  \\[3pt]

\multicolumn{1}{c}{PSM}    &        $^\dagger0.972_{\pm{0.085}}$       &       $^\dagger0.989_{\pm{0.093}}$         & \multicolumn{1}{l}{} &   $^\dagger0.957_{\pm{0.769}}$     &    $^\dagger0.980_{\pm{0.088}}$    &  &   $^\dagger0.940_{\pm{0.060}}$ & $^\dagger0.975_{\pm{0.080}}$      \\[3pt]
\cmidrule{1-9}

\multicolumn{1}{c}{RF}    &        $^\dagger0.987_{\pm{0.0543}}$       &       $^\dagger0.902_{\pm{0.119}}$         & \multicolumn{1}{l}{} &   $^\dagger1.001_{\pm{0.033}}$      &    $^\dagger0.839_{\pm{0.144}}$     &  &   $ ^\dagger1.014_{\pm{0.041}}$ & $^\dagger0.804_{\pm{0.158}}$\\[3pt]
\multicolumn{1}{c}{CF}    &        $^\dagger0.852_{\pm{0.082}}$       &       $^\dagger0.890_{\pm{0.099}}$         & \multicolumn{1}{l}{} &   $^\dagger0.819_{\pm{0.082}}$      &    $^\dagger0.855_{\pm{0.106}}$   &  &   $^\dagger0.786_{\pm{0.076}}$  &  $^\dagger0.826_{\pm{0.102}}$\\[3pt]
\cmidrule{1-9}


\multicolumn{1}{c}{TARNET}    &        $^\dagger0.522_{\pm{0.165}}$       &       $^\dagger0.556_{\pm{0.193}}$         & \multicolumn{1}{l}{} &   $\dagger0.326_{\pm{0.821}}$  &    $\dagger0.363_{\pm{0.101}}$   &     &   $^\dagger0.225_{\pm{0.026}}$& $^\dagger0.260_{\pm{0.041}}$      \\[3pt]
\multicolumn{1}{c}{CFR}    &        $^\dagger0.521_{\pm{0.166}}$       &       $^\dagger0.561_{\pm{0.193}}$         & \multicolumn{1}{l}{} &  $\dagger0.311_{\pm{0.072}}$ &    $\dagger0.348_{\pm{0.861}}$      &  &   $\bf{0.215_{\pm{0.019}}}$ & $\dagger0.258_{\pm{0.036}}	$   	 \\[3pt]
\cmidrule{1-9}

CP~(proposed) &       $\bf 0.307_{\pm{0.125}}$ & $\bf 0.331_{\pm{0.149}}$        & \multicolumn{1}{l}{} &     $\bf 0.249_{\pm{0.063}}$          & $\bf 0.259_{\pm{0.086}}$       &  &   $0.229_{\pm{0.031}} $& ${\bf 0.211_{\pm{0.049}}}$      \\[3pt]

\bottomrule[2pt]
    \end{tabular}
    }
\end{table*}





\subsection{Baselines}
% \subsubsection{Baselines}
We compare the proposed method with several existing supervised ITE estimation approaches.
(i) Linear regression~(Ridge, Lasso) is the ordinary linear regression models with ridge regularization or lasso regularization. 
We consider two variants: one that includes the treatment as a feature~(denoted by `Ridge-$1$' and `Lasso-$1$'),
and the other with two separated models for treatment and control~(denoted by `Ridge-$2$' and `Lasso-$2$').  
% \vspace{1mm}
% \noindent
% {\bf LinearRegression-2:} is separated ordinary linear regression model for treatment and control.
% \vspace{1mm}
% \noindent
(ii) $k$-nearest neighbors~($k$NN) is a matching-based method that predicts the outcomes using nearby instances.
% \vspace{1mm}
% \noindent
(iii) Propensity score matching with logistic regression~(PSM)~\cite{rosenbaum1983central} is a matching-based method using the propensity score estimated by a logistic regression model.
% \vspace{1mm}
% \noindent
% {\bf Random forest~(RF)}~\cite{breiman2001random} and its causal extension called 
% {\bf Causal forest~(CF)}~\cite{wager2018estimation} are also tested.
We also compared the proposed method with tree models such as  (iv) random forest~(RF)~\cite{breiman2001random} and its causal extension called (v) causal forest~(CF)~\cite{wager2018estimation}. 
In CF, trees are trained to predict propensity score and leaves are used to predict treatment effects. 
% \kashima{Include TARNet. I am not sure the difference between TARNet and CFR. Both seem intoduced in Shalit2017.???} \haradap{Actually both are included in shalit2017. TARNet + balancing = CFR. }
% (vi) TARNet~\cite{shalit2017estimating} and (vii) Counterfactual regression~(CFR-Wass)~\cite{shalit2017estimating} are state-of-the-art DNN models based on balanced representations between treatment and control instances. We use the Wasserstein model. 
% \kashima{You can just call CFR (instead of CFRWass)?}
(vi) TARNet~\cite{shalit2017estimating}  is a deep neural network model that has shared layers for representation learning and different layers for outcome prediction for treatment and control instances. 
%See Section~\ref{sec:proposed} for more details.  
(vii) Counterfactual regression~(CFR)~\cite{shalit2017estimating} is a  state-of-the-art deep neural network model based on balanced representations between treatment and control instances. We use the Wasserstein distance. 



\subsection{Results and discussions}
% \kashima{RQs should be the same as the ones we gave in the beginning of this section.}
% \subsubsection{Results and discussions}
% \kashima{Does this section test performance FOR SMALL LABELED DATA?}\haradap{Yes}
%We discuss the experimental results.
% \harada{In this subsection, we mainly discuss how the predictive performance changes according to the size of labeled data and the magnitude of noise. First we explain experimental results varying the size of labeled data and sensitivity to hyper-parameters which control the strength of propagation. Then we analyze the experimental results on noisy outcomes.}
We discuss the performance of the proposed method compared with the baselines by changing the size of labeled datasets, and then investigate the robustness against the label noises. 

We first see the experimental results for different sizes of labeled datasets and sensitivity to the choice of the hyper-parameters that control the strength of label propagation. 
% 全体的な傾向、提案手法良い。NEWSで特に良い
% Tables \ref{table:news} and \ref{table:ihdp} show the PEHE values by different methods for the News dataset and the IHDP dataset, respectively. 
% Overall, our proposed method exhibits the best ITE estimation performance for both labeled and unlabeled data in both of the datasets; the advantage is more significant in the News dataset. 
%\harada{
Tables \ref{table:news}, \ref{table:ihdp}, and \ref{table:synthetic} show the PEHE values by different methods for the News dataset, the IHDP dataset, and the Synthetic dataset, respectively. 
Overall, our proposed method exhibits the best ITE estimation performance for both labeled and unlabeled data in all of the three datasets.
%Among the two semi-synthetic datasets (News and IHDP), the advantage is more significant in the News dataset. 
%}
In general, the performance gain by the proposed method is larger on labeled data than on unlabeled data. 

% NEWSでの話
%On the News dataset, the proposed method achieves the best performance in terms of PEHE for both the labeled and unlabeled instances  settings. 
The News dataset is a relatively high-dimensional dataset represented using a bag of words.
%seems to be relatively easier for defining  appropriate similarity measures. 
The two-model methods such as Ridge-2 and Lasso-2 perform well in spite of their simplicity, and in terms of regularization types, the Lasso-based methods perform relatively better due to the high-dimensional nature of the dataset. 

% IHDPの話 と やや精度が落ちる理由
The proposed method also performs the best in the IHDP dataset; however, the performance gain is rather moderate, as shown by the no statistical significance against CFR~\cite{shalit2017estimating} with the largest $40\%$-labeled data, which is the most powerful baseline method.
The reason for the moderate improvements is probably because of the difficulty in defining appropriate similarities among instances, because the IHDP dataset has various types of features including continuous variables and discrete variables.
%The traditional baselines including Ridge1, Lasso1, kNN matching, and the tree based models do not degrade their predictive accuracy significantly due to the scarcity of labeled data, however their predictive capability still limit their performance even when we increase the labeled data. 
% \kashima{Traditional baselines such as Ridge1, Lasso1, kNN matching, and the tree based models show limited performance; }
The traditional baselines such as Ridge-1, Lasso-1, $k$-NN matching, and the tree-based models show limited performance;
in contrast, the deep learning based methods such as TARNet and CFR demonstrate remarkable performance. 
%However, these methods seem to be vulnerable to the decrease of labeled data.  

% Syntheticの話
%\harada{
The proposed method again achieves the best performance in the Synthetic dataset. 
Since the outcomes are generated by a non-linear function, the linear regression methods are not capable of capturing the non-linear behaviour, and therefore show the limited performance. 
Though TARNet and CFR show the excellent results particularly when using $40$\% of the whole dataset, they suffer from the scarcity of labeled data, and significantly degrade their performance. 
On the other hand, the proposed method is rather robust to the lack of labeled data.
%}



\begin{table}[tbp]
\caption{Investigation of the contributions by the outcome propagation and the ITE propagation in the proposed method.
% The upper table shows the results for the News dataset, and the lower for the IHDP dataset.
The top table shows the results for the News dataset, the middle one for the IHDP, and the bottom one for the Synthetic dataset.
%on the News dataset~(upper)~and the IHDP dataset~(lower). 
The $\lambda_o=0$ and $\lambda_e=0$ indicate the proposed method (CP) with only the ITE propagation and the outcome propagation, respectively,
The $^\dagger$ indicates that our proposed method (CP) performs statistically significantly better than the baselines by the paired $t$-test ($p<0.05)$. 
The bold numbers indicate the best results in terms of the average.}\label{table:proposed_comparison} 
% On IHDP dataset, the outcome regularizations seem to play a more important role than the treatment effect resgularization. On News dataset, when smaller size of labeled data, treatment effect regularization shows slightly better effects in terms of average.}}
\vspace{2mm}
  \centering
\begin{tabular} {ccccccccc }
\toprule[2pt]
        $ \sqrt{\epsilon_{\rm PEHE}} $         & \multicolumn{2}{c}{ \bf{News} 1\%}        &              & \multicolumn{2}{c}{\bf{News} 5\%}& &  \multicolumn{2}{c}{\bf{News} 10\%}      \\[3pt]\cmidrule{2-3}\cmidrule{5-6}\cmidrule{8-9}
Method                      & labeled  & unlabeled & \multicolumn{1}{l}{} & labeled & unlabeled & \multicolumn{1}{l}{} & labeled & unlabeled\\[3pt] \cmidrule{1-9}

CP~($\lambda_o=0$)  &  $\bf 2.812_{\pm{651}}$ & $\bf 2.806_{\pm{0.598}}$  & \multicolumn{1}{l}{} &  $^\dagger2.527_{\pm{0.474}}$  & $\dagger2.531_{\pm{0.523}}$  &     &  $^\dagger2.400_{\pm{0.347}}$ & $^\dagger2.410_{\pm{0.450}}$  \\[3pt]

CP~($\lambda_e=0$) &        $2.879_{\pm{0667}}$       &       $2.885_{\pm{0.609}}$         & \multicolumn{1}{l}{} &     $ 2.351_{\pm{0.450}}$   &   $2.483_{\pm{0.481}}$      &       &  $\bf 1.996_{\pm{0.338}}$ & $2.221_{\pm{0.455}}$   \\[3pt]

CP &        $2.844_{\pm{0.683}}$       &       $2.823_{\pm{0.656}}$        & \multicolumn{1}{l}{} &     $\bf 2.310_{\pm{0.430}}$          & $\bf 2.446_{\pm{0.471}}$       &  &   $2.003_{\pm{0.393}} $& $\bf 2.153_{\pm{0.436}}$      \\[3pt]
\bottomrule[2pt]
\end{tabular}
    \vspace{3mm} \\
    
\centering
\begin{tabular} {ccccccccc }
\toprule[2pt]
        $ \sqrt{\epsilon_{\rm PEHE}} $         & \multicolumn{2}{c}{ \bf{IHDP} 10\%}        &              & \multicolumn{2}{c}{\bf{IHDP} 20\%}& &  \multicolumn{2}{c}{\bf{IHDP} 40\%}      \\[3pt]\cmidrule{2-3}\cmidrule{5-6}\cmidrule{8-9}
Method                      & labeled  & unlabeled & \multicolumn{1}{l}{} & labeled & unlabeled & \multicolumn{1}{l}{} & labeled & unlabeled\\[3pt] \cmidrule{1-9}

CP~($\lambda_o=0$)  &  $^\dagger2.883_{\pm{3.708}}$ & $^\dagger3.004_{\pm{4.071}}$  & \multicolumn{1}{l}{} &  $^\dagger1.972_{\pm{1.930}}$  & $^\dagger2.144_{\pm{2.465}}$  &      &  $1.574_{\pm{1.392}}$  & $1.674_{\pm{1.874}}$      \\[3pt]

CP~($\lambda_e=0$) &        $2.494_{\pm{3.201}}$       &       $2.698_{\pm{3.461}}$         & \multicolumn{1}{l}{} &     $ 1.728_{\pm{2.194}}$   &   $1.977_{\pm{2.450}}$      &      &  $1.344_{\pm{1.383}}$  & $1.585_{\pm{1.923}}$        \\[3pt]

CP &       $\bf 2.427_{\pm{3.189}}$ & $\bf 2.652_{\pm{3.469}}$        & \multicolumn{1}{l}{} &     $\bf 1.686_{\pm{1.838}}$          & $\bf 1.961_{\pm{2.343}}$       &  &   $\bf 1.299_{\pm{1.001}} $& $\bf 1.485_{\pm{1.433}}$      \\[3pt]
\bottomrule[2pt]
    \end{tabular}
        \vspace{3mm} \\
    
% 以下２つ人工データ
\centering
\begin{tabular} {ccccccccc }
\toprule[2pt]
        $ \sqrt{\epsilon_{\rm PEHE}} $         & \multicolumn{2}{c}{ \bf{Synthetic} 10\%}        &              & \multicolumn{2}{c}{\bf{Synthetic} 20\%}& &  \multicolumn{2}{c}{\bf{Synthetic} 40\%}      \\[3pt]\cmidrule{2-3}\cmidrule{5-6}\cmidrule{8-9}
Method                      & labeled  & unlabeled & \multicolumn{1}{l}{} & labeled & unlabeled & \multicolumn{1}{l}{} & labeled & unlabeled\\[3pt] \cmidrule{1-9}

CP~($\lambda_o=0$)  &  $0.334_{\pm{0.112}}$ & $0.367_{\pm{0.124}}$  & \multicolumn{1}{l}{} &  $0.256_{\pm{0.067}}$  & $0.263_{\pm{0.089}}$  &      &  $0.236_{\pm{0.044}}$  & $0.224_{\pm{0.063}}$      \\[3pt]

CP~($\lambda_e=0$) &        $0.320_{\pm{0.110}}$       &       $0.356_{\pm{0.126}}$         & \multicolumn{1}{l}{} &     $ 0.272_{\pm{0.062}}$   &   $0.269_{\pm{0.083}}$      &      &  $0.232_{\pm{0.036}}$  & $0.217_{\pm{0.063}}$        \\[3pt]


CP &       $\bf 0.307_{\pm{0.125}}$ & $\bf 0.331_{\pm{0.149}}$        & \multicolumn{1}{l}{} &     $\bf 0.249_{\pm{0.063}}$          & $\bf 0.259_{\pm{0.086}}$       &  &   $\bf 0.229_{\pm{0.031}} $& $\bf 0.211_{\pm{0.049}}$      \\[3pt]
\bottomrule[2pt]
    \end{tabular}
\end{table}







% ２つのpropagationを比較
Our proposed method has two different propagation terms, the outcome propagation term and the ITE propagation term, as regularizers for semi-supervised learning. 
Table~\ref{table:proposed_comparison} investigates the contributions by the different propagation terms. 
The proposed method using the both propagation terms~(denoted by CP)~shows better results than the one only with the ITE propagation denoted by CP~($\lambda_o=0$); on the other hand, the improvement over the one only with the outcome regularization is marginal.
This observation implies the outcome propagation contributes more to the predictive performance than the ITE propagation.

% hyperparameterへの依存
We also examine the sensitivity of the performance to the regularization hyper-parameters.
% Figure~\ref{fig:hyperparameters} reports the results using $40\%$ and $10\%$ of the whole data as the training data of the News and IHDP datasets, respectively. 
% The proposed method seems rather sensitive to the strength of the regularization terms, particularly on the IHDP dataset, which suggests that the regularization parameters should be carefully tuned using validation datasets in the proposed method. 
%\harada{
Figure~\ref{fig:hyperparameters} reports the results using $10\%$, $20\%$ ,and $40\%$ of the whole data as the training data of the News, IHDP, and Synthetic datasets, respectively. 
The proposed method seems rather sensitive to the strength of the regularization terms, particularly on the IHDP dataset, which suggests that the regularization parameters should be carefully tuned using validation datasets in the proposed method. 
%}
%We also note that the proposed method includes many hyper-parameters which should be fine-tuned. 
In our experimental observations, slight changes in the hyper-parameters sometimes caused significant changes of predictive performance. 
We admit the hyper-parameter sensitivity is one of the current limitations in the proposed method and  efficient tuning of the hyper-parameters should be addressed in future.
% \harada{We compare the proposed method with the state-of-the-art methods varying the magnitude of noise on outcomes. Fig.~\ref{fig:noise_result}~delivers the performance comparison in terms of $\rm \sqrt{\epsilon_{PEHE}}$. Note that the results when $c=1$ correspond to the results in Table~\ref{table:ihdp}~and Table~\ref{table:news}. It is observed that the proposed method is not sensitive when the magnitude of noise is not significant. However, if we add more larger noise, the proposed method fails to propagate outcomes effectively and get more severely affected by the noisy instances than baselines. In both datasets,  we remark that as well as other semi-supervised learning methods, too noisy instances could be crucial as some previous studies have reported~\cite{vahdat2017toward,bui2018neural,du2018robust,liu2012robust}. }

% ノイズを変えて実験
Finally, we compare the proposed method with the state-of-the-art methods by varying the magnitude of noises added to the outcomes. 
Fig~\ref{fig:noise_result}~shows the performance comparison in terms of $\rm \sqrt{\epsilon_{PEHE}}$.
% Note that the results when $c=1$ correspond to the previous results in Tables~\ref{table:news} and \ref{table:ihdp} . 
Note that the results when $c=1$ are the same as those in Tables~\ref{table:news}, \ref{table:ihdp} and \ref{table:synthetic}.
The proposed method stays tolerant of relatively small magnitude of noises; however, with larger label noises, it suffers more from wrongly propagated outcome information than the baselines. 
This is consistent with the previous studies reporting the vulnerability of semi-supervised learning methods against label noises~\cite{vahdat2017toward,bui2018neural,du2018robust,liu2012robust}. 

 \begin{figure}[tb]
%  \begin{minipage}{0.2455\hsize}
\centering
\begin{minipage}{0.43\hsize}
%   \begin{center}
   \includegraphics[width=\linewidth]{images/news_in_reg.pdf}
      \centering
   (a): labeled~(News)
%   \end{center}
 \end{minipage}
 \centering
  \begin{minipage}{0.43\hsize}
%   \begin{center}
%    \includegraphics[width=\linewidth]{images/ihdp_out_reg.pdf}
   \includegraphics[width=\linewidth]{images/news_reg.pdf}
      \centering
   (b): unlabeled~(News)
%   \end{center}
 \end{minipage}\\
 \centering
 \begin{minipage}{0.43\hsize}
%   \begin{center}
   \includegraphics[width=\linewidth]{images/ihdp_in_reg.pdf}
   \centering
   (c): labeled~(IHDP)
%   \end{center}
 \end{minipage}
 \centering
  \begin{minipage}{0.43\hsize}
%   \begin{center}
%    \includegraphics[width=\linewidth]{images/news_out_reg.pdf}
   \includegraphics[width=\linewidth]{images/ihdp_reg.pdf}
      \centering
   (d): unlabeled~(IHDP)
%   \end{center}
 \end{minipage}
  \centering
 \begin{minipage}{0.43\hsize}
%   \begin{center}
   \includegraphics[width=\linewidth]{images/synthetics_in_reg.pdf}
   \centering
   (e): labeled~(Synthetic)
%   \end{center}
 \end{minipage}
 \centering
  \begin{minipage}{0.43\hsize}
%   \begin{center}
%    \includegraphics[width=\linewidth]{images/news_out_reg.pdf}
   \includegraphics[width=\linewidth]{images/synthetics_reg.pdf}
      \centering
   (f): unlabeled~(Synthetic)
%   \end{center}
 \end{minipage}
 \caption{ Sensitivity of the performance to the hyper-parameters.  The colored bars indicate $\sqrt{\epsilon_{\text{PEHE}}}$ for (a)(b) the News dataset, (c)(d) the IHDP dataset  and (e)(f) the Synthetic dataset, when using the largest size of labeled data.
 The deeper-depth colors indicate larger errors.
 It is observed that the proposed method is somewhat sensitive to the choice of the hyper-parameters, especially, the strength of the outcome regularization~($\lambda_o$).  }\label{fig:hyperparameters}
\end{figure}
%
\begin{figure}[t]
 \begin{minipage}{0.33\hsize}
  \begin{center}
   \includegraphics[width=\linewidth]{images/news_noise.pdf}
   (a):  News
  \end{center}
 \end{minipage}
 \begin{minipage}{0.33\hsize}
  \begin{center}
   \includegraphics[width=\linewidth]{images/ihdp_noise.pdf}
   (b):  IHDP
  \end{center}
 \end{minipage}
  \begin{minipage}{0.33\hsize}
  \begin{center}
   \includegraphics[width=\linewidth]{images/synthetic_noise.pdf}
   (c):Synthetic  
  \end{center}
 \end{minipage}
    \caption{Performance comparisons for different levels of noise $c$ added to the labels on (a) News dataset, (b) IHDP dataset, and (c) Synthetic dataset. Note that the results when $c=1$ correspond to the previous results (Tables~\ref{table:news},\ref{table:ihdp}, and \ref{table:synthetic}).}\label{fig:noise_result}
    \end{figure}
    

% レイアウトの問題でここに表を移動
% \begin{table}[tbp]
% \caption{\harada{The performance comparison between the proposed methods on News dataset. $\lambda_o=0$ and $\lambda_e=0$ indicate CP with only treatment effect regularization and outcome regularization. $^\dagger$ indicates that our proposed method (CP) performs statistically significantly better than the baselines by the $t$-test ($p<0.05)$. The bold results indicate the best results in terms of the average.}\label{table:cp_comparison}}
% \centering
% \begin{tabular} {ccccccccc }
% \toprule[2pt]
%         $ \sqrt{\epsilon_{\rm PEHE}} $         & \multicolumn{2}{c}{ \bf{News} 1\%}        &              & \multicolumn{2}{c}{\bf{News} 5\%}& &  \multicolumn{2}{c}{\bf{News} 10\%}      \\[3pt]\cmidrule{2-3}\cmidrule{5-6}\cmidrule{8-9}
% Method                      & labeled  & unlabeled & \multicolumn{1}{l}{} & labeled & unlabeled & \multicolumn{1}{l}{} & labeled & unlabeled\\[3pt] \cmidrule{1-9}

% CP~($\lambda_o=0$)  &  $\bf 2.812_{\pm{651}}$ & $2.806_{\pm{0.598}}$  & \multicolumn{1}{l}{} &  $^\dagger2.527_{\pm{0.474}}$  & $\dagger2.531_{\pm{0.523}}$  &     &  $^\dagger2.400_{\pm{0.347}}$ & $^\dagger2.410_{\pm{0.450}}$  \\[3pt]

% CP~($\lambda_e=0$) &        $2.879_{\pm{0667}}$       &       $2.885_{\pm{0.609}}$         & \multicolumn{1}{l}{} &     $ 2.351_{\pm{0.450}}$   &   $2.483_{\pm{0.481}}$      &       &  $\bf 1.996._{\pm{0.338}}$ & $2.221_{\pm{0.455}}$   \\[3pt]

% CP &        $2.844_{\pm{0.683}}$       &       $\bf 2.823_{\pm{0.656}}$        & \multicolumn{1}{l}{} &     $\bf 2.310_{\pm{0.430}}$          & $\bf 2.446_{\pm{0.471}}$       &  &   $2.003_{\pm{0.393}} $& $\bf 2.153_{\pm{0.436}}$      \\[3pt]
% \bottomrule[2pt]
%     \end{tabular}
% \end{table}




% 実データ→人工データの順番がいいかも?
% 人工と実でそれぞれサブせくしょんにした方が良いか

% \subsection{Experiments on synthetic dataset}
% \harada{Aiming to examine the robustness against the noisy instances, we conduct experiments using synthetic dataset. In this experiments, we vary the magnitude and proportion of noisy outcomes.}

% \subsection{RQ3: Is the proposed model robust to noisy outcomes?}
% \subsubsection{Experimental setting.}

% \harada{
% We generate $1,000$ instances who have $8$ covariates ${\bf x}$ from $\mathcal{N}({\bf 0}^{8\times1},0.5\times(\Sigma+\Sigma^{\top} ))$ where $\Sigma $ $\sim\mathcal{U}((-1,1)^{8\times8})$ and the treatment $t$ from $t\mid\text{Bern}(\sigma({\bf w}_t^{\top}{\bf x}+\epsilon))$ where $\sigma(x)=\frac{1}{1+\exp(-x)}, {\bf w}_t\sim\mathcal{U}((-1,1)^{8\times1}), \epsilon_t\sim\mathcal{N}(0, 0.1)$. We consider the following two types of outcome generative systems:~(i): a non-linear outcome and treatment effect model~(ii): a linear outcome and treatment effect model.
% (i): In the treatment outcome and control outcome is generate as $y^{1}\mid x\sim ({\bf w}^{\top}_{y^{1}}x +c\epsilon_{y})$ and $y^{0}\mid x\sim ({\bf w}^{\top}_{y^{0}}x+c\epsilon_{y})$, where ${\bf w} \sim\mathcal{U}((-1,1)^{8\times1}), \epsilon_y\sim\mathcal{N}({\bf 0}, 0.1)$.  (ii):In the treatment outcome and control outcome is generate as  $y^{1}\mid x\sim (\sin({\bf w}^{\top}_yx) +c\epsilon_{y})$ and $y^{0}\mid x\sim (\cos({\bf w}^{\top}_yx) +c\epsilon_{y})$, where ${\bf w} \sim\mathcal{U}((-1,1)^{8\times1}), \epsilon_y\sim\mathcal{N}({\bf 0}, 0.1)$. $c$ is the parameter which controls the magnitude of noise. 
% Varying the magnitude of noise, we observe how the performance change.
% }
% As we focus on the effect of data scarcity issue, we use $10\%$ of whole data as train data and validation data and use the rest of them as test data.

% \harada{Based on the previous works~\cite{hill2011bayesian,johansson2016learning}, we consider more noisy outcomes for each dataset. In addition to the noiseless potential outcomes, we add noise $c\epsilon$ where $c$ is the magnitude of noise and $\epsilon\sim\mathcal{N}(0, 1)$. Then we vary the magnitude of noise $c$ from $1,3,5,7,9$ and observe how the performances change.
% Note that in both dataset, the outcomes employed in the previous studies~\cite{hill2011bayesian,johansson2016learning,shalit2017estimating} and RQ1  are generated when $c=1$. See the original papers for more detail on the outcome generation~\cite{hill2011bayesian,johansson2016learning}.}

% As we focus on the effect of data scarcity issue, we use $10\%$ of whole data as train data for IHDP dataset and $1\%$ for News dataset.

% \subsubsection{Results and discussions}
% % 人工データ用の説明
% \harada{We compare our method with the state-of-the-art methods varying the magnitude of noisy outcomes. Fig.~\ref{fig:synthetic_result}~demonstrate the performance comparison in terms of $\rm \epsilon_{PEHE}$. It is observed that the proposed method is not sensitive when the magnitude of noise is not significant but if we add more crucial noise, the proposed method seem to fail to propagate outcomes effectively and get more severely affected by the noisy instances. In case~(i), the proposed method yield better performances than baselines when the magnitude is small, however when the magnitude is large~($c=9$), the proposed method is not capable of propagations effectively. In case~(ii), similar phenomena is observed when small magnitude. Standard deviation gets higher than baselines. Though our framework is robust to noisy instances to some extent, we remark that as well as other semi-supervised learning methods, too noisy instances could be more crucial than baselines  some previous studies have noted~\cite{vahdat2017toward,bui2018neural}.
% We do not observe the statistically significant difference between TARNet and CFR.}




% IHDPとNews用の説明
% \harada{We compare the proposed method with the state-of-the-art methods varying the magnitude of noisy outcomes. Fig.~\ref{fig:noise_result}~demonstrate the performance comparison in terms of $\rm \epsilon_{PEHE}$. Note that the results when $c=1$ correspond to the results in Table~\ref{table:ihdp}~and Table~\ref{table:news}. It is observed that the proposed method is not sensitive when the magnitude of noise is not significant. However, if we add more crucial noise, the proposed method seem to fail to propagate outcomes effectively and get more severely affected by the noisy instances than baselines. In both datasets,  we remark that as well as other semi-supervised learning methods, too noisy instances could be more crucial than baselines as some previous studies have reported~\cite{vahdat2017toward,bui2018neural}. }



% It is observed that existing methods which rely only on the labeled data do not seem to work properly . In contrast, the proposed method shows the robustness against the the scarcity of labeled data and effectively. 
% With regard to noisy instances, the proposed method is not sensitive when the magnitude of noise is not significant. However, if we add more crucial noise, the proposed method seem to fail to propagate outcomes effectively and get more severely affected by the noisy instances. In case~(i), the proposed method get degraded given even small proportion of noise~(Fig.~\ref{fig:synthetic_result}(a), (b), (c))~and baselines also failed when large proportion and large noise. We do not observe the statistically significant difference between TARNet and CFR. In case~(ii), the proposed method is not sensitive when the noisy instances are less~(Fig.~\ref{fig:synthetic_result_sin}(a), (b)). However, the proposed method seems to get more affected by noisy instances when the proportion of noise increase and could be crucial if more instances get noisy. CFR-Wass show better results than TARNet. We assume more complex outcome systems of case~(ii) than case~(i)~make such difference.



% \begin{figure}[t]
%  \begin{minipage}{0.33\hsize}
%   \begin{center}
%    \includegraphics[width=\linewidth]{images/synthetic_5.pdf}
%    (a):  $5\%$ of noisy instances
%   \end{center}
%  \end{minipage}
%  \begin{minipage}{0.33\hsize}
%   \begin{center}
%    \includegraphics[width=\linewidth]{images/synthetic_10.pdf}
%    (b):  $10\%$ of noisy instances
%   \end{center}
%  \end{minipage}
%   \begin{minipage}{0.33\hsize}
%   \begin{center}
%    \includegraphics[width=\linewidth]{images/synthetic_20.pdf}
%    (c):  $20\%$ of noisy instances
%   \end{center}
%  \end{minipage}
%    \caption{Performance comparison on synthetic dataset across 10 times experiments in case~(i). Shaded space indicate the standard deviation.}\label{fig:synthetic_result}
%     \begin{minipage}{0.33\hsize}
%   \begin{center}
%    \includegraphics[width=\linewidth]{images/synthetic_sin_5.pdf}
%    (a): $5\%$ of noisy instances
%   \end{center}
% %   \caption{Performance comparison on synthetic dataset across 10 times experiments. Each error bar indicate the standard deviation.}\label{fig:synthetic_result}
%  \end{minipage}
%  \begin{minipage}{0.33\hsize}
%   \begin{center}
%    \includegraphics[width=\linewidth]{images/synthetic_sin_10.pdf}
%    (a): $10\%$ of noisy instances
%   \end{center}
%  \end{minipage}
%   \begin{minipage}{0.33\hsize}
%   \begin{center}
%    \includegraphics[width=\linewidth]{images/synthetic_sin_20.pdf}
%    (a): $20\%$ of noisy instances
%   \end{center}
%  \end{minipage}
%    \caption{Performance comparison on synthetic dataset across 10 times experiments in case~(ii). Shaded space indicate the standard deviation.}\label{fig:synthetic_result_sin}
% \end{figure}


% 人工データ
    
% \begin{figure}[t]
%  \begin{minipage}{0.49\hsize}
%   \begin{center}
%    \includegraphics[width=\linewidth]{images/synthetic_case1.pdf}
%    (a):  Case~(i)
%   \end{center}
%  \end{minipage}
%  \begin{minipage}{0.49\hsize}
%   \begin{center}
%    \includegraphics[width=\linewidth]{images/synthetic_case2.pdf}
%    (b):  Case~(ii)
%   \end{center}
%  \end{minipage}
%    \caption{Performance comparison on synthetic dataset across 10 times experiments in case~(ii). Shaded space indicate the standard deviation.}\label{fig:synthetic_result}
% \end{figure}



% % Please add the following required packages to your document preamble:
% % \usepackage{multirow}
% \begin{table}[]
% \begin{tabular}{cccccc}
%         & \multicolumn{2}{c}{\multirow{2}{*}{\textbf{IHDP 1\%}}} &  & \multicolumn{2}{c}{\multirow{2}{*}{\textbf{IHDP 5\%}}} \\
%         & \multicolumn{2}{c}{}                                   &  & \multicolumn{2}{c}{}                                   \\
% method  & Within-sample             & Without-sample             &  & Within-sample             & Without-sample             \\
% OLS-1   &                           &                            &  &                           &                            \\
% OLS-2   &                           &                            &  &                           &                            \\
% Jaccard &                           &                            &  &                           &                            \\
% katz    &                           &                            &  &                           &                           
% \end{tabular}
% \end{table}
% \clearpage

% \begin{figure}[bt]
%  \begin{minipage}{0.33\hsize}
%   \begin{center}
%    \includegraphics[width=\linewidth]{images/synthetic.pdf}
%   \end{center}
% %   \caption{Performance comparison on synthetic dataset across 10 times experiments. Each error bar indicate the standard deviation.}\label{fig:synthetic_result}
%  \end{minipage}
%  \begin{minipage}{0.33\hsize}
%   \begin{center}
%    \includegraphics[width=\linewidth]{images/synthetic.pdf}
%   \end{center}
%  \end{minipage}
%   \begin{minipage}{0.33\hsize}
%   \begin{center}
%    \includegraphics[width=\linewidth]{images/synthetic.pdf}
%   \end{center}
%  \end{minipage}
%    \caption{Performance comparison on synthetic dataset across 10 times experiments. Each error bar indicate the standard deviation.}\label{fig:synthetic_result}
%     \begin{minipage}{0.33\hsize}
%   \begin{center}
%    \includegraphics[width=\linewidth]{images/synthetic.pdf}
%   \end{center}
% %   \caption{Performance comparison on synthetic dataset across 10 times experiments. Each error bar indicate the standard deviation.}\label{fig:synthetic_result}
%  \end{minipage}
%  \begin{minipage}{0.33\hsize}
%   \begin{center}
%    \includegraphics[width=\linewidth]{images/synthetic.pdf}
%   \end{center}
%  \end{minipage}
%   \begin{minipage}{0.33\hsize}
%   \begin{center}
%    \includegraphics[width=\linewidth]{images/synthetic.pdf}
%   \end{center}
%  \end{minipage}
%    \caption{Performance comparison on synthetic dataset across 10 times experiments. Each error bar indicate the standard deviation.}\label{fig:synthetic_result}
% \end{figure}


% \begin{figure}[t]
%  \begin{minipage}{\hsize}
%   \begin{center}
%    \includegraphics[width=\linewidth,height=200pt]{images/synthetic.pdf}
%   \end{center}
%  \end{minipage}
%    \caption{Performance comparison on synthetic dataset across 10 times experiments. Each error bar indicate the standard deviation.}\label{fig:synthetic_result}
% \end{figure}


